\documentclass{article}

% Recommended packages
\usepackage{microtype}
\usepackage{graphicx}
\usepackage{subfigure}
\usepackage{booktabs}
\usepackage{hyperref}

% ICML style (accepted for coursework / camera-ready style)
\usepackage[accepted]{icml2025}

% Math
\usepackage{amsmath}
\usepackage{amssymb}
\usepackage{mathtools}

% Running title
\icmltitlerunning{Deep Vision-Language Models}

\begin{document}

\twocolumn[
\icmltitle{Deep Vision-Language Models:\\
PA0}

% Anonymous author (safe for coursework)
\begin{icmlauthorlist}
\icmlauthor{Muhammad Aslam Adnan - 25280067}{yyy}
\end{icmlauthorlist}

\icmlaffiliation{yyy}{Department of Electrical Engineering, LUMS}

\icmlcorrespondingauthor{Muhammad Aslam Adnan}{25280067@lums.edu.pk}

% \icmlkeywords{Deep Learning, Transfer Learning, Residual Networks}

\vskip 0.3in
]

\printAffiliationsAndNotice{}

\noindent\textbf{Repository:} \url{https://github.com/adnan821/DVLM-PA0}

\section{Task 1: Inner Workings of ResNet-152}



\begin{abstract}
I investigated representation learning and transfer in modern vision and vision–language models through four controlled experiments. First, I analyzed the inner workings of ResNet-152 under a transfer learning regime on CIFAR-10, evaluating head-only training, the impact of removing residual connections, feature evolution across depth, and fine-tuning trade-offs. Second, I examined Vision Transformers (ViTs) by studying pretrained image classification, patch-level attention visualizations, robustness to patch masking, and linear probing using different token pooling strategies. Third, I trained a variational autoencoder on FashionMNIST, analyzed reconstructions and generations under different latent priors, and investigated posterior collapse through ELBO decomposition and KL annealing. Finally, I evaluated CLIP on STL-10 using zero-shot classification with multiple prompting strategies, analyzed the modality gap between image and text embeddings, and assessed post-hoc alignment using Procrustes analysis. Across tasks, the results highlighted the effectiveness of transfer learning, the importance of architectural inductive biases, the trade-offs between representation quality and optimization stability, and the role of embedding geometry in multimodal models.
\end{abstract}
\subsection{Introduction}
Deep residual networks enable the training of very deep models by alleviating optimization difficulties through skip connections. In this work, I empirically analyzed ResNet-152 on a small target dataset to understand feature transferability, gradient flow, and representation quality. I follow a controlled set of experiments aligned with standard transfer learning practice.

\subsection{Experimental Setup}
\textbf{Model.} I use a ResNet-152 model pretrained on ImageNet. The final classification layer is replaced to predict 10 CIFAR-10 classes.

\textbf{Training Protocol.} Unless otherwise stated, training is performed for 5 epochs using standard data augmentation. Training and validation accuracy are reported.

\textbf{Baselines.} I consider (i) head-only training with a frozen backbone, (ii) modified networks with selected residual connections disabled, and (iii) transfer learning variants including pretrained versus random initialization and partial versus full fine-tuning.
\subsection{Baseline: Head-Only Transfer Learning}
I froze the backbone and train only the classification head.

\textbf{Results.} Validation accuracy increases rapidly and saturates within a few epochs, reaching approximately 96.3\% by epoch five.

\textbf{Discussion.} Training ResNet-152 from scratch on small datasets is unnecessary and impractical due to severe overfitting risk given the parameter count and excessive computational requirements. Freezing most of the network demonstrates that early and mid-level features learned on ImageNet—such as edges, textures, and object parts—are highly transferable, requiring only a lightweight task-specific classifier.

\subsection{Residual Connections in Practice}
I disable skip connections in selected residual blocks and retrain the classification head.

\textbf{Results.} Removing residual connections degrades validation accuracy to approximately 85.7\% and slows convergence relative to the baseline.

\textbf{Analysis.} Skip connections provide identity paths that improve gradient flow, reduce vanishing gradients, and enable stable optimization in very deep networks. When residuals are removed, convergence becomes slower and final performance drops, highlighting the central role of residual learning.

\subsection{Feature Hierarchies and Representations}
I extract features from early, middle, and late layers of the network and visualize them using dimensionality reduction techniques such as t-SNE and UMAP.

\textbf{Observations.}
Early layers show weak class separability and are dominated by low-level cues such as edges and colors. Middle layers begin to cluster semantically related classes, while late layers exhibit clear class separation, indicating high-level, task-aligned representations. This progression confirms that discriminative structure emerges gradually with network depth.

\subsection{Transfer Learning and Generalization}
I compared ImageNet-pretrained models with randomly initialized models and evaluate different fine-tuning strategies.

\textbf{Results.}
Random initialization with limited fine-tuning achieves approximately 69--70\% validation accuracy, whereas ImageNet-pretrained models with full-backbone fine-tuning reach approximately 86.4\%.

\textbf{Trade-offs.}
Fine-tuning only the final block provides a favorable compute--accuracy compromise, while full fine-tuning yields the best accuracy at higher computational cost. Early layers are the most transferable, while deeper layers benefit from task-specific adaptation.

\subsection{Optional Analysis}
\textbf{t-SNE vs. UMAP.}
UMAP preserves global structure more effectively and produces more compact clusters, while t-SNE emphasizes local separation.

\textbf{Class Confusions.}
Misclassifications typically occur between visually similar CIFAR-10 classes such as cats and dogs, reflecting shared high-level features.

\textbf{Depth Comparison.}
ResNet-152 features are more linearly separable than those from ResNet-18, though gains diminish relative to increased computational cost.

\subsection{Conclusion}
Our experiments demonstrate that transfer learning with frozen backbones is highly effective on small datasets, residual connections are essential for optimization in very deep networks, and feature representations become increasingly discriminative with depth. ImageNet pretraining provides substantial benefits, and selective fine-tuning offers an effective balance between accuracy and efficiency.

\newpage
\section{Task 2: Understanding Vision Transformers (ViT)}

\subsection{Using a Pre-trained ViT for Image Classification}
I used the ImageNet-pretrained model \texttt{google/vit-base-patch16-224} and its corresponding image processor to preprocess images to the required $224\times 224$ resolution. I evaluated the model on two test images.

\textbf{Results.} The top-1 predictions were:
\begin{itemize}
    \item \textbf{Image 0:} \emph{golden retriever} with probability $p=0.885$.
    \item \textbf{Image 1:} \emph{tiger cat} with probability $p=0.441$.
\end{itemize}
Qualitatively, both predictions appear reasonable for the chosen images. The first prediction is highly confident, while the second has noticeably lower confidence, suggesting greater ambiguity or more competing classes.

% Optional: include the two input images (upload to Overleaf as figs/input0.png and figs/input1.png)
\begin{figure}[t]
\vskip 0.1in
\centering
\subfigure[Input image 0]{\includegraphics[width=0.45\columnwidth]{Screenshot 2026-02-01 at 5.53.50 AM}}
\hfill
\subfigure[Input image 1]{\includegraphics[width=0.45\columnwidth]{Screenshot 2026-02-01 at 5.54.23 AM}}
\caption{Two evaluation images used for ViT top-1 prediction testing.}
\label{fig:vit-inputs}
\vskip -0.1in
\end{figure}

\subsection{Visualizing Patch Attention}
A key advantage of Vision Transformers is attention-based interpretability. I configured the model to output attention weights (\texttt{output\_attentions=True}) and extracted the final-layer attention tensor of shape $(B, H, L, L)$, where $B$ is batch size, $H$ is the number of heads, and $L = N+1$ tokens (patch tokens plus the special \texttt{[CLS]} token).

I aggregated attention across heads by averaging over $H$, then extracted the \texttt{[CLS]} token’s attention to all patch tokens, yielding a vector of length $N=196$ (for $14\times 14$ patches with $16\times16$ patch size). This vector was reshaped into a $14\times14$ attention map, upsampled to image resolution, and overlaid as a semi-transparent heatmap on the original image.

% Optional: include attention overlays (upload to Overleaf as figs/attn0.png and figs/attn1.png)
\begin{figure}[t]
\vskip 0.1in
\centering
\subfigure[Attention overlay (image 0)]{\includegraphics[width=0.45\columnwidth]{1.png}}
\hfill
\subfigure[Attention overlay (image 1)]{\includegraphics[width=0.45\columnwidth]{2.png}}
\caption{Final-layer \texttt{[CLS]} attention over patches, averaged across heads, upsampled and overlaid on the input images.}
\label{fig:vit-attn}
\vskip -0.1in
\end{figure}

\subsection{Attention Map Analysis and Interpretability}
From the attention overlays (Fig.~\ref{fig:vit-attn}), the highlighted regions generally align with the most salient parts of the object for the predicted class, indicating that the ViT is using meaningful evidence rather than purely background cues. Compared with CNN interpretability methods such as CAM/Grad-CAM (conceptually), ViT attention offers a more direct mechanism: attention weights are natively produced by the architecture, allowing patch-level attributions without requiring gradient-based post-hoc methods.

Different attention heads can focus on different structures (e.g., object boundaries vs.\ texture vs.\ context). Aggregating across heads provides a stable summary, but individual heads may show specialization; this can be assessed by visualizing head-specific overlays and comparing their focus patterns.

\subsection{Masking Patches at Inference: Robustness to Missing Information}
I evaluated robustness by masking a fraction of input patches at inference and measuring whether the masked prediction matches the original top-1 class (agreement).

\textbf{Results (from notebook outputs).}
\begin{itemize}
    \item \textbf{Random masking agreement:} $0.600$
    \item \textbf{Structured (center) masking agreement:} $1.000$
\end{itemize}

\textbf{Discussion.} Random masking can remove scattered informative patches and disrupt global evidence aggregation, leading to more frequent prediction changes. Structured masking (center masking) can sometimes be less harmful if the most informative evidence lies outside the masked region for that specific image, or if sufficient redundant cues remain. Overall, ViTs may be reasonably robust to missing patches because self-attention can integrate distributed context; however, robustness depends strongly on which regions are masked and how much redundancy exists in the image.

\subsection{Linear Probes: CLS Token vs Mean Pooling of Patch Tokens}
I compared linear probes trained on (i) the \texttt{[CLS]} embedding versus (ii) the mean of patch token embeddings.

\textbf{Results (validation accuracy over 5 epochs).}
\begin{itemize}
    \item Epoch 1: CLS $0.946$, Mean $0.944$
    \item Epoch 2: CLS $0.944$, Mean $0.938$
    \item Epoch 3: CLS $0.952$, Mean $0.942$
    \item Epoch 4: CLS $0.956$, Mean $0.948$
    \item Epoch 5: CLS $0.956$, Mean $0.948$
\end{itemize}

\textbf{Analysis.} CLS pooling performed slightly better and more consistently than mean pooling in these experiments. This is expected because the CLS token is explicitly optimized (during standard ViT classification pretraining) to aggregate global information for prediction. Mean pooling can be competitive when pretraining objectives encourage uniform information distribution across patch tokens, but under standard supervised ImageNet pretraining, CLS is often the strongest single representation for linear probing.

\subsection{Summary}
Task 2 demonstrates that a pretrained ViT can produce reasonable top-1 predictions on test images, provides built-in patch-level attention visualizations for interpretability, exhibits non-trivial robustness behavior under patch masking (with clear differences between random and structured masking), and slightly favors CLS-based pooling over mean pooling for linear probes in our setting.


% ================= TASK 3 =================
\newpage
\section{Task 3: Training Variational Autoencoders}

\subsection{Train the VAE}

\textbf{Dataset.} I trained on \texttt{FashionMNIST} using \texttt{torchvision}. The dataset was split into train/val/test sizes of 55,000 / 5,000 / 10,000.

\textbf{Model.} I used the provided VAE from \texttt{architecture.py}. The model had 1,691,369 parameters and latent dimension $d_z=20$. Training ran on \texttt{cuda}.

\textbf{Objective.} I optimized the negative ELBO using a reconstruction term (MSE) and KL divergence:
\[
\mathcal{L}(x) \;=\; \underbrace{\|x-\hat{x}\|_2^2}_{\text{reconstruction (MSE)}} \;+\; \beta \cdot 
\underbrace{D_{\mathrm{KL}}\!\left(q_\phi(z\mid x)\,\|\,p(z)\right)}_{\text{KL}}
\]
where $p(z)=\mathcal{N}(0,I)$ and $\beta$ is annealed during training.

\textbf{Training.} I trained for 10 epochs with KL warm-up (linear $\beta$ increase). I recorded train/validation losses and KL (per-epoch averages from notebook output).

\begin{table}[t]
\caption{VAE training/validation losses with KL warm-up.}
\label{tab:vae-losses}
\vskip 0.15in
\centering
\begin{small}
\begin{tabular}{c c c c c}
\toprule
Epoch & $\beta$ & Train Loss & Val Loss & Train KL \\
\midrule
1  & 0.20 & 27.0593 & 18.7150 & 15.6323 \\
2  & 0.40 & 20.4838 & 19.7742 & 13.0539 \\
3  & 0.60 & 21.6254 & 21.4406 & 11.0579 \\
4  & 0.80 & 23.1232 & 23.2129 & 9.6399 \\
5  & 1.00 & 24.6483 & 24.6752 & 8.5944 \\
6  & 1.00 & 24.3725 & 24.6375 & 8.4992 \\
7  & 1.00 & 24.1628 & 24.5551 & 8.4513 \\
8  & 1.00 & 24.0344 & 24.2026 & 8.4653 \\
9  & 1.00 & 23.8666 & 24.0697 & 8.4324 \\
10 & 1.00 & 23.8020 & 23.9665 & 8.4323 \\
\bottomrule
\end{tabular}
\end{small}
\vskip -0.1in
\end{table}

\subsection{Visualize Reconstructions and Generations}

\textbf{Reconstructions.} I passed test images through the encoder-decoder and visualized original (top) vs reconstruction (bottom). Reconstructions captured the rough shape and category-specific structure (e.g., shoes vs shirts) but remained blurry, which is typical for VAEs with pixel-wise MSE.

\begin{figure}[t]
\vskip 0.1in
\centering
\includegraphics[width=\columnwidth]{vae_recon.png}
\caption{Reconstructions on FashionMNIST: originals (top row) and reconstructions (bottom row).}
\label{fig:vae-recon}
\vskip -0.1in
\end{figure}

\textbf{Generations (Normal prior).} I sampled $z\sim\mathcal{N}(0,I)$ and decoded to image space. Generated samples were recognizable but often smooth/averaged, reflecting uncertainty across plausible pixels.

\textbf{Generations (Laplacian prior).} I also sampled from a Laplace prior (same latent dimension) and decoded. Qualitatively, Laplace sampling produced similar coarse shapes but with slightly different texture/contrast patterns in some samples, consistent with heavier-tailed latent draws.

\begin{figure}[t]
\vskip 0.1in
\centering
\subfigure[Normal prior samples]{\includegraphics[width=0.47\columnwidth]{vae_gen_normal.png}}
\hfill
\subfigure[Laplace prior samples]{\includegraphics[width=0.47\columnwidth]{vae_gen_laplace.png}}
\caption{Generated FashionMNIST samples from different latent priors.}
\label{fig:vae-gen}
\vskip -0.1in
\end{figure}

\subsection{Posterior Collapse Investigation}

Posterior collapse occurs when the encoder ignores the input and $q_\phi(z\mid x)$ matches the prior too closely, causing the KL term to shrink toward zero while reconstructions remain dominated by the decoder.

\textbf{ELBO components.} Fig.~\ref{fig:vae-elbo} shows reconstruction loss and KL over training. The KL term decreases from $\approx 15.63$ (epoch 1) to $\approx 8.43$ (epoch 10), but it does \emph{not} vanish; this suggests the encoder is still using $z$ meaningfully (i.e., no complete collapse). However, the generated samples remain visually smooth and somewhat repetitive, which can be a symptom of \emph{partial} collapse or an overly smooth decoder likelihood induced by MSE.

\begin{figure}[t]
\vskip 0.1in
\centering
\includegraphics[width=\columnwidth]{vae_elbo_components.png}
\caption{ELBO components during training (reconstruction vs KL).}
\label{fig:vae-elbo}
\vskip -0.1in
\end{figure}

\textbf{Posterior statistics.} For a minibatch, the latent means had $\mu$ mean $\approx 0.0059$ and $\mu$ std $\approx 0.5653$ (shape $64\times 20$), indicating non-trivial variation across inputs rather than a degenerate constant posterior.

\textbf{Why collapse happens.} Collapse is more likely when:
(i) the decoder is powerful enough to model $p(x\mid z)$ without needing $z$,
(ii) the KL penalty is enforced strongly early in training, pushing $q_\phi(z\mid x)$ toward $p(z)$ before the encoder learns informative codes.

\subsection{Mitigating Posterior Collapse}

\textbf{Strategy.} I used \textbf{KL annealing / warm-up}: start with a small $\beta$ and gradually increase it to 1.0. This gives the encoder time to learn informative latent representations before the KL regularizer is fully enforced.

\textbf{Effect.} With annealing, KL begins higher and decreases smoothly rather than collapsing immediately. The reconstructions remain qualitatively meaningful, and generations are more diverse than the ``single-blob'' failure mode. Fig.~\ref{fig:vae-anneal} visualizes reconstruction and KL dynamics under the annealing regime implemented in the notebook.

\begin{figure}[t]
\vskip 0.1in
\centering
\includegraphics[width=\columnwidth]{vae_elbo_annealing.png}
\caption{ELBO components under KL annealing (mitigation experiment).}
\label{fig:vae-anneal}
\vskip -0.1in
\end{figure}

\textbf{Why it works.} Early training focuses on reconstruction, encouraging $z$ to encode input-dependent information. As $\beta$ increases, the KL term regularizes the learned posterior toward the prior while preserving useful structure, reducing the chance that the encoder collapses to an input-independent distribution.

\subsection{Summary}

I trained a VAE on FashionMNIST using an MSE+KL ELBO objective, visualized reconstructions and generations (Normal vs Laplace priors), analyzed ELBO components for signs of posterior collapse, and applied KL warm-up as a mitigation strategy. Overall, KL remained non-zero, suggesting the model avoided complete collapse, while annealing provided more stable training dynamics and more meaningful latent usage.

% ------------------------
% NOTE FOR COMPILATION:
% Please upload/rename the following images into your Overleaf project:
%   vae_recon.png            (from notebook: /mnt/data/task3_figs/cell11_img0.png)
%   vae_gen_normal.png       (from notebook: /mnt/data/task3_figs/cell12_img1.png)
%   vae_gen_laplace.png      (from notebook: /mnt/data/task3_figs/cell12_img3.png)
%   vae_elbo_components.png  (from notebook: /mnt/data/task3_figs/cell13_img0.png OR cell10_img0.png)
%   vae_elbo_annealing.png   (from notebook: /mnt/data/task3_figs/cell15_img0.png)
% ------------------------


% ================= TASK 4 =================
\newpage
\section{Task 4: Modality Gap in CLIP}

Contrastive Language--Image Pretraining (CLIP) jointly trains a vision encoder and a text encoder to map images and natural language descriptions into a shared embedding space. This enables zero-shot classification by comparing image embeddings with text embeddings derived from human-readable labels, eliminating the need for a task-specific classifier. In this task, I evaluate CLIP on the STL-10 dataset, analyze the modality gap between image and text embeddings, and investigate whether simple post-hoc alignment can bridge this gap.

\subsection{Zero-Shot Classification on STL-10}

\textbf{Dataset and Model.}  
I used the STL-10 dataset from \texttt{torchvision}, evaluating on the standard test split. The pretrained CLIP ViT-B/32 model was loaded from OpenAI’s official implementation and used in evaluation mode without fine-tuning.

\textbf{Prompting Strategies.}  
I evaluated zero-shot classification accuracy using three prompting strategies:
\begin{itemize}
    \item \textbf{Plain labels:} raw class names (e.g., ``cat'').
    \item \textbf{Template prompts:} ``a photo of a \{class\}''.
    \item \textbf{Descriptive prompts:} more detailed natural language descriptions.
\end{itemize}

For each strategy, text embeddings were computed once and reused for the entire test set. Classification was performed via cosine similarity between image and text embeddings.

\textbf{Results.}  
Zero-shot accuracies on the STL-10 test set are reported in Table~\ref{tab:clip-zeroshot}.

\begin{table}[t]
\caption{Zero-shot classification accuracy on STL-10 using different CLIP prompting strategies.}
\label{tab:clip-zeroshot}
\vskip 0.15in
\centering
\begin{small}
\begin{tabular}{l c}
\toprule
Prompting Strategy & Accuracy \\
\midrule
Plain labels & 0.9625 \\
``a photo of \{class\}'' & \textbf{0.9736} \\
Descriptive prompts & 0.9728 \\
\bottomrule
\end{tabular}
\end{small}
\vskip -0.1in
\end{table}

Template-based prompts achieved the highest accuracy, confirming that CLIP is sensitive to prompt formulation and benefits from minimal linguistic context.

\subsection{Exploring the Modality Gap}

\textbf{Embedding Extraction.}  
I randomly sampled 100 STL-10 images and extracted their corresponding image embeddings using the CLIP vision encoder. For each image, the text embedding of its ground-truth class was extracted using the same prompt template. All embeddings were 512-dimensional.

\textbf{Dimensionality Reduction.}  
I applied UMAP to project both image and text embeddings into a two-dimensional space. Embeddings were L2-normalized prior to projection.

\textbf{Observations.}  
The resulting visualizations show a clear modality gap: image and text embeddings occupy partially distinct regions of the embedding space. Text embeddings are more compact and tightly clustered, while image embeddings exhibit greater spread. Although embeddings corresponding to the same class are relatively closer across modalities, the distributions do not fully overlap.

Normalization reduces scale differences but does not eliminate the modality gap. Despite this separation, CLIP performs well because zero-shot classification depends on relative cosine similarity rather than absolute embedding overlap.

\subsection{Bridging the Modality Gap via Procrustes Alignment}

\textbf{Alignment Method.}  
To reduce the modality gap, I applied an orthogonal Procrustes alignment. Given paired image embeddings $X$ and text embeddings $Y$, I solve:
\[
\min_{R} \|XR - Y\|_F,
\]
where $R$ is an orthogonal matrix. I used the library implementation \texttt{scipy.linalg.orthogonal\_procrustes}.

\textbf{Procedure.}  
Paired embeddings from the same 100 samples were used to learn the rotation matrix $R \in \mathbb{R}^{512 \times 512}$. Image embeddings were rotated using $R$, while text embeddings remained fixed.

\textbf{Effect on Embeddings.}  
After alignment, dimensionality-reduced visualizations show that image embeddings move closer to their corresponding text embeddings, and the modality gap is visibly reduced. However, the distributions do not collapse completely, indicating that a single global linear transform is insufficient for perfect alignment.

\textbf{Classification After Alignment.}  
Zero-shot classification accuracy using the aligned embeddings is reported below:

\begin{center}
\begin{tabular}{l c}
\toprule
Setting & Accuracy \\
\midrule
Before alignment & \textbf{0.9736} \\
After Procrustes alignment & 0.9711 \\
\bottomrule
\end{tabular}
\end{center}

Alignment slightly reduces accuracy, suggesting that CLIP’s original embedding geometry is already well-optimized for cosine similarity, and post-hoc linear alignment can disrupt fine-grained semantic relationships.

\subsection{Discussion}

These experiments show that CLIP exhibits a measurable modality gap between image and text embeddings, even though they lie in a shared embedding space. Prompt engineering has a substantial impact on zero-shot performance, often more than explicit embedding alignment. While Procrustes alignment improves visual overlap between modalities, it does not improve classification accuracy, highlighting that CLIP’s success relies on relative similarity structure rather than full modality collapse.

\subsection{Summary}

In this task, I evaluated CLIP’s zero-shot performance on STL-10, analyzed the modality gap via low-dimensional projections, and applied orthogonal Procrustes alignment to bridge this gap. Despite visible modality separation, CLIP achieves high accuracy through robust similarity geometry. Simple linear alignment reduces the visual gap but does not improve—and slightly degrades—classification performance.


\end{document}
